\documentclass[12pt]{ximera}
\title{Exam 1 (2021)}
\begin{document}
\begin{abstract}
An archived exam on Chapter 1: majority thresholds, Plurality and Plurality-with-Elimination winners, fairness criteria, Condorcet candidates, a modified Borda count, and Pairwise Comparisons point totals.
\end{abstract}
\maketitle

\section*{Exam 1 (2021)}

\begin{problem}
Problems 1--4 refer to the following preference schedule for an election with four
candidates.
\begin{center}
\begin{tabular}{c|ccc}
Number of voters & 8 & 6 & 5\\\hline
1st & B & D & C\\
2nd & D & B & A\\
3rd & A & C & D\\
4th & C & A & B
\end{tabular}
\end{center}

\begin{prompt}
\textbf{1.} To obtain a majority of the votes cast, a candidate needs

\begin{multipleChoice}
\choice{$9$ votes}
\choice{$19$ votes}
\choice{more votes than any other candidate}
\choice[correct]{$10$ votes}
\end{multipleChoice}

\begin{feedback}
There are $8+6+5=19$ voters, and a majority is more than half: $10$.
\end{feedback}
\end{prompt}

\begin{prompt}
\textbf{2.} Which candidate wins under the Plurality method?

\begin{multipleChoice}
\choice{A}
\choice[correct]{B}
\choice{C}
\choice{D}
\choice{There is a tie}
\end{multipleChoice}

\begin{feedback}
First-place votes: B $8$, D $6$, C $5$, A $0$.
\end{feedback}
\end{prompt}

\begin{prompt}
\textbf{3.} Which candidate wins under Plurality-with-Elimination?

\begin{multipleChoice}
\choice{A}
\choice{B}
\choice{C}
\choice[correct]{D}
\choice{The method is not required / does not apply}
\end{multipleChoice}

\begin{feedback}
Eliminate A ($0$), then C ($5$). C's voters rank D next (A is gone), giving D $11$ and
B $8$. D wins.
\end{feedback}
\end{prompt}

\begin{prompt}
\textbf{4.} In this election, a Condorcet candidate

\begin{multipleChoice}
\choice{is A}
\choice{is B}
\choice{is C}
\choice[correct]{is D}
\choice{does not exist}
\end{multipleChoice}

\begin{feedback}
D beats A $14$--$5$, B $11$--$8$, and C $14$--$5$.
\end{feedback}
\end{prompt}
\end{problem}

\begin{problem}
\textbf{5.} True or false: for any election, the Condorcet criterion is never violated by
the Pairwise Comparisons method.

\begin{multipleChoice}
\choice[correct]{True}
\choice{False}
\end{multipleChoice}

\begin{feedback}
True. A Condorcet candidate wins every comparison and so has the most points.
\end{feedback}
\end{problem}

\begin{problem}
\textbf{6.} True or false: the Majority criterion is violated by the Plurality method.

\begin{multipleChoice}
\choice{True}
\choice[correct]{False}
\end{multipleChoice}

\begin{feedback}
False. A candidate with a majority of first-place votes necessarily has the most
first-place votes, so always wins by Plurality.
\end{feedback}
\end{problem}

\begin{problem}
Consider the following preference schedule.
\begin{center}
\begin{tabular}{c|cccc}
Number of voters & 5 & 4 & 2 & 1\\\hline
1st & C & A & B & D\\
2nd & D & C & A & C\\
3rd & A & B & D & B\\
4th & B & D & C & A
\end{tabular}
\end{center}

\begin{prompt}
\textbf{7.} Use Pairwise Comparisons to find a complete ranking.

A: $\answer{2}$ \quad B: $\answer{0.5}$ \quad C: $\answer{2.5}$ \quad D: $\answer{1}$
points

\smallskip
1st $\answer{C}$, 2nd $\answer{A}$, 3rd $\answer{D}$, 4th $\answer{B}$

\begin{feedback}
\[
\begin{array}{lll@{\qquad}lll}
\text{A vs B} & 9\text{--}3 & \text{A} & \text{B vs C} & 2\text{--}10 & \text{C}\\
\text{A vs C} & 6\text{--}6 & \text{tie} & \text{B vs D} & 6\text{--}6 & \text{tie}\\
\text{A vs D} & 6\text{--}6 & \text{tie} & \text{C vs D} & 9\text{--}3 & \text{C}
\end{array}
\]
\end{feedback}
\end{prompt}

\end{problem}

\begin{problem}
An election with four candidates (A, B, C, D) and $60$ voters used a modified Borda
count: $6$ points for first place, $4$ for second, $2$ for third, and $1$ for fourth. A
had $140$ points, B had $200$, and C had $320$.

\begin{prompt}
\textbf{(a)} How many points did D have? $\answer{120}$

\begin{feedback}
Each ballot awards $6+4+2+1=13$ points, so the total is $60\cdot13=780$. D has
$780-(140+200+320)=120$.
\end{feedback}
\end{prompt}

\begin{prompt}
\textbf{(b)} Give the ranking of the candidates.

1st $\answer{C}$, 2nd $\answer{B}$, 3rd $\answer{A}$, 4th $\answer{D}$

\begin{feedback}
C $320$, B $200$, A $140$, D $120$.
\end{feedback}
\end{prompt}
\end{problem}

\begin{problem}
An election with seven candidates (A--G) used the Pairwise Comparisons method. Candidate
A lost $1$ comparison; B lost $2$ and tied $1$; C lost $2$ and tied $2$; D lost $3$ and
tied $1$; F lost $5$; and G lost $6$.

\begin{prompt}
\textbf{(a)} How many points did E get? $\answer{6}$

\begin{hint}
Each comparison awards exactly $1$ point in total. How many comparisons are there?
\end{hint}

\begin{feedback}
Each candidate is in $6$ comparisons, and a win is worth $1$, a tie $\frac12$:
\[
\text{A } 5,\quad \text{B } 3.5,\quad \text{C } 3,\quad \text{D } 2.5,\quad
\text{F } 1,\quad \text{G } 0,
\]
which total $15$. There are $\frac{7\cdot6}{2}=21$ comparisons, so E has
$21-15=6$ points: E won every comparison.
\end{feedback}
\end{prompt}

\begin{prompt}
\textbf{(b)} Give a complete ranking of the candidates.

1st $\answer{E}$, 2nd $\answer{A}$, 3rd $\answer{B}$, 4th $\answer{C}$, 5th $\answer{D}$,
6th $\answer{F}$, 7th $\answer{G}$

\begin{feedback}
E $6$, A $5$, B $3.5$, C $3$, D $2.5$, F $1$, G $0$.
\end{feedback}
\end{prompt}
\end{problem}

\end{document}
